\documentclass[14pt,a4paper]{extarticle}
\usepackage[T2A]{fontenc}
\usepackage[utf8]{inputenc}
\usepackage[english,russian]{babel}
\usepackage{lmodern}
\usepackage{geometry}
\usepackage{array}
\usepackage{tabularx}
\usepackage{longtable}
\usepackage{ltablex}
\usepackage{booktabs}
\usepackage{hyperref}
\usepackage{indentfirst}
\usepackage{microtype}
\usepackage{graphicx}
\usepackage{setspace}
\usepackage{listings}
\usepackage{graphicx}
\usepackage{caption}
\usepackage{float}
\usepackage{tocloft}
\usepackage{titlesec}
\usepackage{newtxtext,newtxmath}

\titleformat{\section}        % форматируем \section
  {\normalfont\bfseries}      % стиль: обычный + жирный
  {\thesection.}              % номер секции с точкой
  {2mm}                       % расстояние между номером и текстом
  {}                           % перед заголовком можно добавить что-то
\titlespacing*{\section}{1.5cm}{*3}{*1}  % {отступ слева}{отступ сверху}{отступ снизу}

\captionsetup{justification=centering}
\renewcommand{\ttdefault}{cmtt}
\keepXColumns
\setlength{\parindent}{3em}
\newgeometry{left=2.5cm,right=2.5cm,top=3cm,bottom=2.5cm}

\renewcommand{\contentsname}{\hspace*{4em}{\smallСодержание}}
\setlength{\cftbeforesecskip}{0.5em}
\renewcommand{\cftsecleader}{\cftdotfill{\cftdotsep}}
\renewcommand{\cftsecaftersnum}{.\hspace{2mm}}

\title{Отчёт по дисциплине "Основы предпринимательства"}
\author{Выполнил студент группы М-РИТ-25-1, Станиславчук Сергей}

\begin{document}

\setstretch{1.5}
\noindent
\begin{minipage}[t]{0.01\textwidth}
    \includegraphics{~/programmer/++Programmer/5_1/LSTU.jpg}
\end{minipage}%
\begin{minipage}[c]{1.1\linewidth}
    \centering
    {\small
    МИНИСТЕРСТВО НАУКИ И ВЫСШЕГО ОБРАЗОВАНИЯ РФ \\[1em]
    ФЕДЕРАЛЬНОЕ ГОСУДАРСТВЕННОЕ БЮДЖЕТНОЕ \\
    ОБРАЗОВАТЕЛЬНОЕ УЧРЕЖДЕНИЕ ВЫСШЕГО ОБРАЗОВАНИЯ \\
    «ЛИПЕЦКИЙ ГОСУДАРСТВЕННЫЙ ТЕХНИЧЕСКИЙ УНИВЕРСИТЕТ»
    }
\end{minipage}%

\vspace{3em}

\noindent
\hspace{1em}%
\begin{minipage}[t]{0.8\textwidth}
\textbf{Институт} \hspace{2em} компьютерных наук \\
\textbf{Кафедра} \hspace{2em} автоматизированных систем управления
\end{minipage}


\vspace{6em}

\begin{center}
\textbf{ПРАКТИЧЕСКАЯ РАБОТА №4} \\
по дисциплине \\ <<Основы предпринимательства>>
\end{center}

\vspace{6em}

\noindent
\begin{tabular*}{\textwidth}{@{\extracolsep{\fill}} l c r }
    Студент М-РИТ-25-1 & \makebox[4cm]{\hrulefill} & Станиславчук С.М. \\[2em]
Руководитель & \\
Канд. юр. наук, доцент & \makebox[4cm]{\hrulefill} & Моисеева И.И. \\
\end{tabular*}

\vfill
\begin{center}
    Липецк, 2025
\end{center}

\thispagestyle{empty}

\newpage
\tableofcontents
\newpage

\section{Задание 1}
\begin{sloppypar}
Организация применяет упрощенную систему налогообложения. В
качестве объекта налогообложения выбраны «доходы, уменьшенные на величину
расходов». В 1 квартале текущего года были осуществлены следующие операции:
— Реализовано продукции на сумму — 620 тыс. руб.
— Приобретено и оплачено сырье для производства продукции на сумму 350
тыс.руб., НДС — 63 тыс, руб.
- Произведена оплата за аренду производственного помещения в сумме 150 тыс. руб.
и НДС — 27 тыс, руб.
- Затраты на оплату труда составили — 102 тыс, руб.
- Перечислены взносы во внебюджетные социальные фонды (30,2 \%) в сумме 34,7
тыс.руб.
Определите налоговую базу и сумму налога по упрощенной системе
налогообложения за 1 квартал текущего года

Решение.

    \texttt{620 - 350 - 63 - 150 - 27 - 102 - 34.7 = -106.7} - убыток => единый налог не начисляется

\end{sloppypar}

\section{Задание 2}
ООО «Колос» находится на территории Липецкой области и занимается
производством кукурузы и ее первичной переработкой. Также организация предоставляет
услуги в области гостиничного бизнеса и общественного питания для населения в летнее
время. По итогам календарного года общество получило доход в размере 22,5 млн руб., в
том числе от предоставленных услуг в размере 4,6 млн руб. Определите, имеет ли право
ООО «Колос» перейти на уплату единого сельскохозяйственного налога.

Решение.

Для перехода на ЕСХН организция должна выполнять два ключевых условия:
\newline1. Не менее 70\% доходов - от сельхозпроизводства
\newline2. Организация должна быть сельхозтоваропроизводителем
1) \texttt{22.5 - 4.6 = 17.9 млн руб} - доход от сельхоздеятельности
\texttt{17.9 / 22.5 = 79.6\%} - доля сельхоздохода
Требование 1 выполнено.
2)ООО «Колос» производит кукурузу и занимается её первичной переработкой. Это
полностью попадает под сельхозпроизводство.

Ответ: имеет право

\section{Задание 3}
Задание 3. АО «Родина» является сельскохозяйственным товаропроизводителем, с
2005 г. перешло на уплату единого сельскохозяйственного налога. В соответствии с
налоговым учетом по результатам прошлого года получило доходы от реализации 600 500
тыс.руб., из них от сельскохозяйственной продукции — 570 500 тыс.руб.;
внереализационные доходы — 230 тыс. руб. Расходы, исчисленные в соответствии с
требованиями налогового учета для сельскохозяйственных товаропроизводителей,
составили за налоговый период 580 700 тыс, руб. и внереализационные расходы — 120
тыс, руб.
Определите единый сельскохозяйственный налог для АО «Родина» за прошлый год.
Назовите основные элементы налога, порядок перехода на уплату единого
сельскохозяйственного налога, отличие специального режима от общей системы
налогообложения

1) Доходы от реализации: 600 500 тыс. руб.
Из них сельхозпродукция: 570 500 тыс. руб.

Доля сельхоздоходов:
\texttt{570 500 / 600 500 = 95\%} => условие >70\% выполнено.

Переход на ЕСХН правомерен.

2) \texttt{600 500 + 230 = 600 730 тыс. руб.} - доходы

\texttt{580 700 + 120 = 580 820 тыс. руб.} - расходы

\texttt{600 730 – 580 820 = 19 910 тыс. руб.} - налоговая база

\texttt{19 910 × 6\% = 1 194,6 тыс. руб.} - налог (ставка 6\%)

ЕСХН:
Объект налогообложения: доходы минус расходы.

Налоговая база: денежное выражение объекта.

Ставка: 6%.

Налоговый период: календарный год.

Отчетный период: полугодие.

Плательщики: организации и ИП, у которых >=70\% доходов — от сельхозпроизводства.

Порядок перехода:
\newlineПодать заявление в налоговую до 31 декабря предшествующего года (для перехода
со следующего года).

Вновь созданные — в течение 30 дней с регистрации.

Условие: не менее 70\% доходов от собственной сельхозпродукции.

Отличия от общей системы:
\newlineНет НДС, налога на прибыль и налога на имущество (кроме отдельных объектов).

Простая налоговая база: доходы минус расходы.

Ставка фиксированная — 6\%.

Меньше отчетности и проще учет.

На ОСНО — налог на прибыль, НДС, имущество + сложный учет и отчетность.

Ответ: ЕСХН к уплате АО «Родина»: 1 194,6 тыс. руб.

\section{Задание 4}

Организация-производитель табачных изделий приобрела 1,5 т табака по
цене 177 руб. за кг (без НДС). Из табака произведены сигареты в количестве 60 коробок по
пять блоков, в каждом блоке по 10 пачек, в пачке по 20 сигарет. На каждой пачке указано,
что розничная цена не может превышать 160 руб.
Рассчитайте сумму акциза к уплате.

Решение.

Специфическая часть - я не знаю наверняка, предположу 2804 руб./1000 сигарет.

\texttt{2 804 * (60 000 / 1 000) = 2 804 * 60 = 168 240 руб.} - специфическая часть
\texttt{160 * 3 000 = 480 0000 руб.} - акцизная база
\texttt{480 0000 * 0.16} - адвалорная часть
\texttt{168 240 + 76 800 = 245 040 руб.} - общий акциз

Ответ: 245 040 рублей (специфическая ставка: 2 804 руб за 1 000 сигарет)

\section{Задание 5}
Работнице организации, не состоящей в зарегистрированном браке (вдове),
имеющей ребёнка в возрасте 12 лет ежемесячно за период январь—май начислялась
зарплата 25тыс.руб. В повторный брак работница вступила 20 февраля. Найдите
сумму налога на доходы физических лиц за период январь—май

Решение.

1) Январь-февраль - двойной вычет 2 800 руб.

\texttt{ 25 000 - 2 800 = 22 200 руб. } - налогооблагаемая база\newline
\texttt{ 22 200 * 13\% = 2 886 руб. } - НДФЛ за месяц\newline
\texttt{ 2 886 * 2 = 5 772 руб. } - за два месяца\newline

2) Март-май (обычный вычет 1 400 руб.)

\texttt{ 25 000 - 1 400 = 23 600 руб. } - налогоооблагаемая база\newline
\texttt{ 23 600 * 13\% = 3 068 руб. } - НДФЛ за месяц\newline
\texttt{ 3 068 * 3 = 9 204 руб. } - за три месяца\newline

Ответ:
5 772 + 9 204 = 14 976 рублей

\section{Задание 6}
Организация за 2019 г. реализовала товаров на 2 800 тыс. руб. (без НДС).
Расходы на производство товаров составили 1 900 тыс. руб., в том числе расходы на
приобретение призов во время рекламной кампании —31 тыс. руб. Кроме того, было
реализовано два транспортных средства:
— автомобиль — 30 июня 2019 г. по цене 180 тыс. руб. (без НДС), первоначальная
стоимость — 270 тыс. руб., сумма амортизационных отчислений — 54 тыс. руб., срок
эксплуатации — два года, срок полезного использования — 10 лет;
— трактор — 25 ноября 2019 г. по цене 170 тыс. руб. (без НДС), первоначальная
стоимость — 380 тыс. руб., сумма амортизационных отчислений —237 500 руб., срок
эксплуатации — пять лет, срок полезного использования — восемь лет.
Определите налог на прибыль, подлежащий уплате в бюджет

Решение.

1) Прибыль от реализации товаров
Из расходов нужно исключить призы, т.к. расходы на рекламные призы учитываются
в пределах 1\% от выручки = 28 тыс.руб.

1 900 - (31 - 28) = 1897 тыс. руб. - уточненные расходы
2 800 - 1 897 = 903 тыс. руб. - прибыль от товаров

2) Реализация автомобиля
180 - 216 = -36 тыс.руб.

3) Реализация трактора
170 - 142.5 = 27.5 тыс.руб.

4) Совокупная прибыль для налогообложения
903 - 36 + 27.5 = 178.9 тыс. руб.

Ответ: 178 900 рублей

\section{Задание 7}
Предприятие за год реализовало 200 тыс. продукции по цене 7,08 руб. (с
НДС), расходы предприятия на з/п сотрудникам составили по 10 тыс. руб. в месяц (всего З
чел.) (без учета взносов в фонды). В этом году предприятие принимало делегацию и
потратило на банкет 30 тыс. руб. и 10 тыс. руб. на экскурсию.
Материальные расходы предприятия составили 2 руб. за единицу (сумма без НДС),
расходы за коммунальные услуги —240 тыс. руб., в том числе НДС.
Определите налог на прибыль организации.

Решение.

7.08 / 1.2 = 5.90 руб. - цена без НДС
5.9 * 200 000 = 1 180 000 руб. - выручка без НДС
10 000 руб. * 12 * 3 = 360 000 руб. - зарплата
200 000 * 2 = 400 000 руб. - мат. расходы
240 000 / 1.2 = 200 000 руб. - коммунальные услуги
4\% * 360 000 = 14 400 руб. - представлительские расходы (не более 4\% от
расходов на оплату труда)
360 000 + 400 000 + 200 000 + 14 000 = 974 400 руб. - общая сумма расходов для налога
1 180 000 - 974 000 = 205 600 руб. - налоговая база
205 600 * 0.2 - 41 120 руб. - налог на прибыль

Ответ: 41 120 рублей

\section{Задание 8}
Организация занимается грузоперевозками. На балансе числятся 5
автомобилей МАЗ 5335 мощностью двигателя 180 л.с. и 4 ЗИЛ 133 мощностью двигателя
120 л.с. Организация 10 августа продала один автомобиль ЗИЛ 133.
Рассчитайте налог, который должна уплатить организация за налоговый период.

Решение.

Посчитал на сайте www.nalog.gov.ru для Липецкой области при 12 полных месяцах владения:
180 л.с. = 9 000 руб / 12 мес. 2025 года,
120 л.с. = 4 800 руб / 12 мес. 2025 года

9 000 * 5 = 45 000 руб. - МАЗ
4 800 * 3 = 14 400 руб.

4 800 / 12 * 8 = 3 000 руб. - налог за проданный ЗИЛ

14 400 + 3 200 = 17 600 руб - ЗИЛ

45 000 + 17 600 = 62 600 руб. - общий налог

Ответ: 62 600 рублей.

\section{Задание 9}
На балансе предприятия числится имущество, остаточная стоимость
составляет:
на 01.01.2020 г. —2 345 000 руб.,
на 0 1.02.2020 г. —2 294 700 руб.,
на 01.03.2020 г. —2 175,300 руб.,
на 01.04.2020 г. —3187600руб.
Определите среднюю стоимость имущества и авансовый платеж по налогу на
имущество организаций

Решение.

1. 2 345 000 + 2 294 700 = 4 639 700
2. 4 639 700 + 2 175 300 = 6 815 000
3. 6 815 000 + 3 187 600 = 10 002 600
10 002 600 / 4 - 2 500 650 руб.
Ставка налога на имущество организаций в РФ не выше 2.2\%
2 500 650 * 0.0222 * 0.25 - 13 753 руб.

\section{Задание 10}
Гражданин имеет в собственности 2-комнатную квартиру кадастровой
стоимостью 2 850 тыс. р. и площадью 52 м2.
Каков размер налога на имущество физических лиц?

Решение.

20/52 * 2 850 000 = 1 096 154 руб.
2 850 000 - 1 096 154 = 1 753 846 руб.

Ставка налога 0.1\% (до 10 млн. руб.)
1 753 846 * 0.001 = 1 754 руб.

Ответ: 1 754 руб.

\section{Тестирование}
\begin{enumerate}
    \item а) индивидуальный, обязательный, безвозмездный взнос денежных средств.
    \item а) б)
    \item а) б) в) д)
    \item б)
    \item а) г)
    \item а)
    \item б)
    \item а)
    \item в)
    \item а)
    \item в)
    \item а)
    \item а
    \item в)
    \item а) г)
\end{enumerate}

\end{document}
